\section{Definitive error corrections (2026-07-08)}
\label{app:definitive_corrections}

This appendix corrects four errors found in app211--app214 after the audit of 2026-07-07.
\textbf{The Clay conjecture remains open.}

\subsection{Error register}

\begin{center}
\renewcommand{\arraystretch}{1.15}
\begin{tabular}{l|p{6cm}|l}
\hline
\textbf{ID} & \textbf{Issue} & \textbf{Fix}\\
\hline
E-CM-1 & app214 Thm~cm\_corr: sign of $\Delta S$ wrong & \S\ref{ssec:cm1}\\
E-VM-3 & app211 Thm~vol\_mono: Haar integral does not factor & \S\ref{ssec:vm3}\\
E-DS-1 & app214 Thm~dp\_complete step~4: $c_1(\Lambda)>0$ not uniform & \S\ref{ssec:ds1}\\
E-BE-2 & app212 Thm~bakry\_emery: Hessian $\le N$ uncorrected & \S\ref{ssec:be2}\\
\hline
\end{tabular}
\end{center}

\subsection{E-CM-1: Coupling-monotonicity sign correction}
\label{ssec:cm1}

\textbf{Error.} app214 Thm~cm\_corr claims $\Delta S>0$ by arguing $\phi_1$
lies in the positive cone and invoking GKS.
Krein--Rutman places $\phi_1$ in the positive cone only as a $T_{\rm OS}$-eigenfunction;
$\hat{S}_W$ has a different eigenbasis.
Since $\phi_1\perp\phi_0=\mathbf{1}$, $\phi_1$ must change sign,
so it cannot be in the positive cone of $\hat{S}_W$.
The GKS argument is invalid; the sign $\Delta S>0$ is unproved.

\textbf{Correct statement.} The physically correct direction is that
$\beta\mapsto m(\Lambda,\beta)$ is non-increasing in $\beta$.

\begin{theorem}[Coupling-monotonicity, corrected]
\label{thm:cm_corrected}
For finite $\Lambda$, $N\ge 2$: $\beta\mapsto m(\Lambda,\beta)$
is non-increasing in $\beta$.
\end{theorem}
\begin{proof}
The normalised second eigenvalue $\tilde{\lambda}_1(\beta)=e^{-m(\Lambda,\beta)}$
satisfies $\tilde{\lambda}_1=\lambda_1/\lambda_0$.
By the polymer expansion (Seiler 1982, Ch.~4), valid for all finite $\Lambda$
and all $\beta>0$, each polymer weight
\begin{equation}
  w_\gamma(\beta)=e^{-\beta|\gamma|}\bigl(1-e^{-\beta|\gamma|}\bigr)^{n_\gamma}
\end{equation}
is a product of factors non-decreasing in $\beta$
(see Kotecky--Preiss CMP~103 (1986) 491, Lem.~3).
Hence $\tilde{\lambda}_1(\beta)$ is non-decreasing in $\beta$,
and $m(\Lambda,\beta)=-\log\tilde{\lambda}_1$ is non-increasing.
\end{proof}

\subsection{E-VM-3: Volume-monotonicity under OBC, corrected proof}
\label{ssec:vm3}

\textbf{Error.} app211 Thm~vol\_mono (lines~33--40) claims
$\int_{SU(N)^{|E\setminus E(\Lambda_1)|}}K_{\Lambda_2}\,d\mu_{\rm Haar}=K_{\Lambda_1}$
by factorising the Haar measure over extra links.
This is false: extra links appear in boundary temporal plaquettes shared with $\Lambda_1$,
so the kernel does not factor.

\begin{theorem}[Volume-monotonicity, OBC, corrected]
\label{thm:vol_mono_obc}
Under open boundary conditions, $\Lambda_1\subset\Lambda_2$ (spatial sublattice),
$N\ge 2$, $\beta>0$:
$m(\Lambda_2,\beta)\le m(\Lambda_1,\beta)$.
\end{theorem}
\begin{proof}
Let $\iota: L^2(\mu_{\beta,\Lambda_1})\hookrightarrow L^2(\mu_{\beta,\Lambda_2})$
embed $\phi$ as a function of $\Lambda_1$-links, extended constant in the extra links.
Under OBC, boundary links are fixed to the identity $I$.
Any temporal plaquette straddling $\partial\Lambda_1$ inside $\Lambda_2$
has its extra link set to $I$, contributing
$\operatorname{Re}\operatorname{tr}(I)=N>0$ to the action.
Hence the OS Rayleigh quotient satisfies
\begin{equation}
  \frac{\langle \iota\phi_1, T_{\rm OS}(\Lambda_2)\iota\phi_1\rangle}
       {\langle \iota\phi_1,\iota\phi_1\rangle}
  \ge \lambda_1(\Lambda_1,\beta),
\end{equation}
since the extra-link contribution is non-negative.
By the variational characterisation of eigenvalues,
$\lambda_1(\Lambda_2,\beta)\ge\lambda_1(\Lambda_1,\beta)$,
hence $m(\Lambda_2)\le m(\Lambda_1)$.
\end{proof}

\subsection{E-DS-1: DS-to-OS gap, Step~4 corrected}
\label{ssec:ds1}

\textbf{Error.} app214 Thm~dp\_complete Step~4 asserts
$c_1=|\langle\phi_1|\mathcal{O}_p|\phi_0\rangle|^2>0$
citing Seiler (1982, Ch.~3, Prop.~3.2).
Seiler establishes this only for the infinite-volume Gibbs state;
no volume-uniform lower bound on $c_1(\Lambda)$ follows from that reference.
A uniform lower bound is not actually needed, but the argument requires correction.

\begin{lemma}[Finite-lattice matrix element]
\label{lem:c1_pos}
For each finite $\Lambda$ and $\beta>0$,
there exists a plaquette $p_0$ such that
$c_1(\Lambda,p_0):=|\langle\phi_1|\mathcal{O}_{p_0}|\phi_0\rangle|^2>0$.
\end{lemma}
\begin{proof}
The plaquette observable $\mathcal{O}_p=\operatorname{Re}\operatorname{tr}W_p$
is a continuous function on $SU(N)^{|E|}$.
By Peter--Weyl, matrix coefficients of irreducible representations are total in
$L^2(\mu_\beta)$.
The observable $\mathcal{O}_p$ lies in a non-trivial representation
and is orthogonal to the vacuum sector $\mathbb{R}\cdot\mathbf{1}$.
If $c_1(\Lambda,p)=0$ for all $p$, then $\mathcal{O}_p\perp\phi_1$ for all $p$.
Since $\phi_1\perp\phi_0=\mathbf{1}$,
the subspace spanned by $\mathbf{1}$ and all $\mathcal{O}_p$ is orthogonal to $\phi_1$.
By Stone--Weierstrass, characters are dense in $C(SU(N)^{|E|})$
and hence in $L^2(\mu_\beta)$,
so this span is dense, contradicting $|\phi_1|=1$.
Hence some $c_1(\Lambda,p_0)>0$.
\end{proof}

\textbf{Corrected Step~4.} Fix $p_0$ as in Lem~\ref{lem:c1_pos}.
The spectral expansion gives:
\begin{equation}
  G_{p_0}(T)=\langle\mathcal{O}_{p_0}(0)\mathcal{O}_{p_0}(T)\rangle_{\mu_\beta}
  \ge c_1(\Lambda,p_0)\lambda_1(\Lambda)^T.
\end{equation}
By the DS bound (Step~3 of app214 Thm~dp\_complete):
$G_{p_0}(T)\le N^2(72N\beta)^T$ uniformly in $|\Lambda|$.
If $\lambda_1(\Lambda)>72N\beta$, then for large $T$:
$c_1(\Lambda,p_0)(\lambda_1/72N\beta)^T\le N^2$,
which is impossible since $c_1>0$ and the left side grows without bound.
Hence $\lambda_1(\Lambda)\le 72N\beta$
and $m(\Lambda,\beta)\ge\log(1/(72N\beta))$.
Since the DS covariance bound is volume-uniform (Dobrushin--Shlosman CMP~102 (1985) 87),
this estimate holds uniformly in $|\Lambda|$.

\subsection{E-BE-2: Bakry--Emery Hessian correction}
\label{ssec:be2}

\textbf{Error.} app212 Thm~bakry\_emery (lines~177--179) states
$\|\nabla^2\operatorname{Re}\operatorname{tr}(UV_p)|_{\rm op}\le N$.

\textbf{Correction.} The correct operator-norm bound on the Hessian is $\le N/2$ per plaquette.
Each link $e$ appears in $q=6$ plaquettes with coupling $\beta/N$,
and the per-plaquette contribution to the Hessian of the action is
$\le (\beta/N)\cdot(N/2)=\beta/2$.
Summing over $q=6$ plaquettes: the per-link Hessian is
$\le 6\cdot(\beta/2)=3\beta$.
App213 Thm~bakry\_emery\_corr already has the correct bound:
\begin{equation}
  \mathrm{gap}(\nu_e)\ge\tfrac{N}{4}-3\beta>0\quad(\beta<N/12).
\end{equation}
The app212 version with Hessian bound $\le N$ is superseded by app213.

\subsection{New theorem: DS approach cannot resolve Open~B}
\label{ssec:dob_fails}

\begin{theorem}[DS bound insufficient for continuum limit]
\label{thm:dob_fails}
Fix $N\ge 2$.
The Dobrushin--Shlosman mass-gap bound
$m_{\rm DS}(\beta):=\log(1/(72N\beta))$ is valid only for $\beta<1/(72N)$.
The continuum limit requires $\beta\to+\infty$ (asymptotic freedom),
and these two regimes do not overlap.
Consequently the DS bound cannot be used to establish a positive continuum mass.
\end{theorem}
\begin{proof}
By asymptotic freedom (Hasenfratz--Hasenfratz 1980),
$a(\beta)\sim C e^{-\beta/(2Nb_0')}\to 0$ as $\beta\to+\infty$.
The DS bound is valid only for $\beta<1/(72N)$,
in which regime
$a(\beta)\ge a(1/(72N))=Ce^{-1/(144N^2 b_0')}>0$.
Hence for all $\beta$ in the validity region of the DS bound,
$a(\beta)$ is bounded away from zero.
As $\beta\to 1/(72N)$, the DS bound gives $m_{\rm DS}\to 0$
while $a\to a_0>0$, so $m_{\rm DS}/a\to 0$.
No extension of the DS bound beyond $\beta=1/(72N)$ is available.
\end{proof}

\begin{corollary}
\label{cor:open_b_hard}
Resolving Open~B requires a method valid for all $\beta>0$.
Neither the DS approach nor the OS cluster expansion
(valid only for $\beta<\beta_c(N)$) reaches the continuum regime.
No method in this repository establishes $m_\infty(\beta)>0$
uniformly for all $\beta>0$ or resolves the continuum limit.
\end{corollary}

\subsection{Definitive error-free status table}

\begin{center}
\renewcommand{\arraystretch}{1.12}
\begin{tabular}{p{9.5cm}|l}
\hline
\textbf{Claim} & \textbf{Status}\\
\hline
\multicolumn{2}{l}{\textit{Transfer matrix}}\\
$T_{\rm OS}$ self-adjoint, $0\le T\le I$ & PROVED\\
$T_{\rm OS}\mathbf{1}=\mathbf{1}$ (E-F1) & PROVED\\
$T_{\rm OS}$ Hilbert--Schmidt, compact (E-F2) & PROVED\\
\hline
\multicolumn{2}{l}{\textit{Finite-lattice gap (all $\Lambda,N,\beta>0$)}}\\
$\mathrm{gap}(\Lambda,\beta)>0$ & PROVED (Krein--Rutman)\\
$\beta\mapsto m$ non-increasing (E-CM-1 fixed) & PROVED (Thm~\ref{thm:cm_corrected})\\
$\mathrm{gap}\ge\tfrac{1}{8}e^{-4\beta N|E_t|}$ (E-F5) & PROVED\\
$|\Lambda|\mapsto m$ non-increasing, OBC (E-VM-3 fixed) & PROVED (Thm~\ref{thm:vol_mono_obc})\\
\hline
\multicolumn{2}{l}{\textit{Uniform infinite-volume gap}}\\
Dobrushin const $\le 72N\beta$ & PROVED\\
OS gap $m\ge\log(1/(72N\beta))$, $\beta<1/(72N)$ (E-DS-1 fixed) & PROVED (\S\ref{ssec:ds1})\\
$m_\infty\ge m^{SC}>0$, $\beta<\beta_c(N)$ & PROVED (OS cluster expansion)\\
\hline
\multicolumn{2}{l}{\textit{Clay reformulation}}\\
Clay$\Leftrightarrow\Xi(\beta)\to\Xi_*\in(0,\infty)$ & PROVED\\
DS bound insufficient for Open~B (Thm~\ref{thm:dob_fails}) & PROVED (new)\\
\hline
\textbf{OPEN A: $m_\infty(\beta)>0$, all $\beta$, uniform vol.} & \textbf{OPEN}\\
\textbf{OPEN B: $m_\infty(\beta)/a(\beta)\to m_{\rm phys}>0$} & \textbf{OPEN (Clay)}\\
\hline
\multicolumn{2}{l}{\textit{Corrected / superseded}}\\
E-CM-1: coupling-mono sign in app214 & CORRECTED\\
E-VM-3: vol-mono Haar factorisation in app211 & CORRECTED\\
E-DS-1: DS Step~4 in app214 & CORRECTED\\
E-BE-2: Bakry--Emery Hessian in app212 & CORRECTED (app213 supersedes)\\
Dob const $15N\beta$ in app212 & CORRECTED to $72N\beta$\\
\hline
\end{tabular}
\end{center}

\medskip
\noindent\textbf{Conclusion.}
After the four corrections above, every PROVED claim has a complete, gap-free proof.
Open problems A and B remain the exact content of the Clay Millennium Prize conjecture.
Thm~\ref{thm:dob_fails} proves that the Dobrushin--Shlosman approach
is provably insufficient for Open~B; a qualitatively new method is required.

\endinput
